% Corsin Week 5, transcript/week-05.md
\chapter{Week 5 --- Optimization, Lagrange Multipliers \& the Hessian Test}

\exercisesheet{5}

\textit{Statements quoted verbatim from \texttt{exercises/Ex5\_Analysis2\_eng.pdf} (assigned 16
March 2026, due 23 March 2026). Attribution: Prof. Joaquim Serra, D-MATH, ETH Z\"urich.}

\begin{notation}[Corsin's recommendations]
\begin{center}
\begin{tabular}{lll}
\textbf{Problem} & \textbf{Priority} & \textbf{Corsin's note} \\
\hline
5.1 & important & ``Recall the linearity of the differential and that $C^1 \Rightarrow$ \\
    & & differentiable'' \\
5.2 & important & ``Compare with the example from the script and recall that continuous \\
    & & functions on compact sets attain their extrema.'' \\
5.3 & important & ``There is no difference between $x \to 0$ and $|x| \to 0$.'' \\
5.4 & semi-important & ``Use the definition of directional derivative, \\
    & & $\partial_v f(x) = \left.\frac{\partial}{\partial t}\right\|_{t=0} f(x+tv)$. Solve a \\
    & & few of these, but not all.'' \\
5.5 & optional & --- \\
5.6 & optional & ``Interesting, but irrelevant for the course'' \\
5.7 & optional & ``Interesting for physicists, will be covered in classical mechanics.'' \\
\end{tabular}
\end{center}
\end{notation}

\begin{notation}
Only the exercises tagged \textbf{important} above are reproduced below. The full sheet
(5.1--5.7, including 5.4--5.7) is transcribed in \texttt{content/exercise-sheets/week-05.tex}.
\end{notation}

\begin{exercise}[5.1 --- Derivatives \textnormal{(important)}]
\label{ex:5.1}
\begin{enumerate}[label=\textbf{(\alph*)}]
  \item Consider $f(x,y) := xy^2e^{-x^2-y^2}$ for $(x,y) \in \mathbb{R}^2$. Find all the critical
    points of $f$, that is all the points where the gradient of $f$ vanishes. (The point will be
    given if and only if all the numerical values you find are correct\ldots\ so check your
    computations twice.)
  \item We want to find a function $g \in C^1(\mathbb{R}^2)$ with the following directional
    derivative:
    $$\partial_v g(x,y) = 2\cos(x^2y)xv_2 + \cos(x^2y)v_1^2y \quad \text{for all } (x,y) \text{ and } (v_1,v_2) \in \mathbb{R}^2.$$
    Give an explicit example of such a function $g$ or prove that a function with these
    properties cannot exist.
\end{enumerate}
\end{exercise}

\begin{remark}
Corsin's hint --- \textit{recall the linearity of the differential} --- is the whole of
\cref{ex:5.1}(b): a directional derivative must be \textbf{linear} in $v$, and $v_1^2$ is not.
\end{remark}

\begin{exercise}[5.2 --- Lagrange multipliers \textnormal{(important)}]
\label{ex:5.2}
Consider the ball $U := \{(x,y,z) : x^2+y^2+z^2 < 1\}$ and the function
$$f(x,y,z) := 1 - x^2 - y^2 + 4x.$$
\begin{enumerate}[label=\textbf{\arabic*.}]
  \item Motivate rigorously why $f$ attains its absolute maximum and minimum in $\overline{U}$.
  \item Find all the critical points of $f$ which lie in the interior of $U$. Say whether
    $\sup_U f$ (or $\inf_U f$) is attained at some point in $U$.
  \item Say whether $\max_{\partial U} f$ and $\min_{\partial U}$ exist.
  \item With the method of Lagrange multipliers, find the possible critical points for the
    constrained problem $\max\{f(x,y,z) : (x,y,z) \in \partial U\}$, and same for min.
  \item Among all the points you have found, clarify at which points the maximum/minimum of $f$
    is attained.
\end{enumerate}
\end{exercise}

\begin{exercise}[5.3 --- Multiple choice (Landau notation) \textnormal{(important)}]
\label{ex:5.3}
Mark all and only the statements which are true.
\begin{enumerate}[label=\textbf{(\alph*)}]
  \item As $|x| \to 0$, if $f(x) = x_2x_1^2 + O(x_1^2)$, then $f(x) = O(|x|^2)$.
  \item As $x \to 0$, if $f(x) = x_2x_1^2 + O(x_1^3)$, then $f(x) = O(x_1^3)$.
  \item As $|x| \to 0$, if $f(x) = x_2x_1 + O(|x|^2)$, then $f$ is differentiable at $x = 0$.
  \item As $x \to 0$, if $f(x) = x_2x_1 + O(|x|^2)$, then $f$ is twice differentiable around
    $x = 0$.
  \item As $|x| \to 0$, if $f(x) = x_2x_1 + O(|x|^3)$, then $f$ is twice differentiable around
    $x = 0$ and $\partial_{11}f(0) = 0$.
\end{enumerate}
\textit{Official hint:} revise \textit{Corollary 10.45}.
\end{exercise}

\begin{importantremark}
Unlike Weeks 2--4, Corsin marks no Monday / Friday split in this week's notes; the material runs
continuously from optimization through Lagrange multipliers to the Hessian test.
\end{importantremark}

\section{Optimization}
\label{sec:optimization}

\begin{definition}
\begin{itemize}
  \item \newterm{Local minimum} (\germanterm{lokales Minimum}) of
    $f : U \subseteq \mathbb{R}^n \to \mathbb{R}$: $f$ has a local minimum at $x_0 \in U$ if
    $$\exists\,\varepsilon > 0 \text{ such that } \forall x \in B_\varepsilon(x_0)\cap U : f(x) \geq f(x_0).$$
  \item \textbf{Local maxima} are defined analogously with $f(x) \leq f(x_0)$.
  \item \textbf{Strict} min/max $\iff f(x) > f(x_0)$ / $f(x) < f(x_0)$.
  \item A \newterm{critical point} (\germanterm{kritischer Punkt}) of $f \in C^1(U,\mathbb{R})$,
    $U \subseteq \mathbb{R}^n$ open, is a point $x_0 \in U$ where $\nabla f(x_0) = 0$.
\end{itemize}
\end{definition}

\begin{proposition}[local extremum $\Rightarrow$ critical point]
\label{prop:extremum_implies_critical}
If $f \in C^1(U,\mathbb{R})$ and $x_0 \in \operatorname{int}(U)$ is a local extremum, then
$\nabla f(x_0) = 0$.
\end{proposition}

\begin{remark}[intuition]
$$\langle\nabla f(x_0), v\rangle = 0 \iff \nabla f(x_0) \perp v,$$
$$\text{and this for all } v \in \mathbb{R}^2 \iff \nabla f(x_0) = 0.$$
\end{remark}

\begin{figurestub}{FIG-W05-01}
A blue goblet-shaped surface $f(x,y)$ standing on an $xy$-grid ($y$-axis to the upper left,
$x$-axis to the right); a purple closed level-set curve is drawn on the grid directly under the
surface's rim. A green dot labelled ``minimum'' sits on the grid at the level curve; a dotted
orange line rises from it to the curve, and a horizontal orange arrow labelled ``gradient''
points left from that dot. A red arrow labelled ``steepest ascent'' starts at the green dot and
climbs up the outside of the surface, following its curvature, to the rim. (Reference asset:
\texttt{Toby Lane/geogebra/gradient\_contour.ggb}.)
\end{figurestub}

\section{Lagrange multipliers}
\label{sec:lagrange_multipliers}

Suppose we want to find the hottest place on the surface of the earth. In theory, temperature is
defined everywhere, giving some function
$$T \in C^\infty(\mathbb{R}^3, (0,\infty)) \qquad \text{(absolute zero is not attainable)}$$
If we model the earth as $S^2 := \{x \in \mathbb{R}^3 : |x| = 1\}$, then we want to find the local
maxima of $T|_{S^2}$, which \textbf{are not necessarily critical points of $T$}!

\begin{figurestub}{FIG-W05-02}
A blue sphere $S^2$ drawn in 3D with a dashed equator ellipse. Near the top, a point $q$ has a
dotted red arrow $\nabla T(q)$ pointing straight outward along the radius (perpendicular to the
sphere), labelled ``local extremum of $T|_{S^2}$''. Lower down, on the front of the sphere, a
point $p$ has a dotted orange arrow $\nabla T(p)$ pointing obliquely up-and-out (not radial ---
visibly tilted toward the surface), labelled ``has a component pointing along the tangent, not a
local extremum of $T|_{S^2}$, since $T$ increases along a direction on $S^2$.''
\end{figurestub}

So $T|_{S^2}$ has a local extremum if it is \textbf{perpendicular to the surface} $S^2$. Note
that $S^2$ is the \newterm{level set} (\germanterm{Niveaumenge}) $g(x) = 0$ for
$g(x) = |x|^2 - 1$, or $S^2 = g^{-1}\{0\}$.

\begin{remark}
$\nabla g(x) = \dfrac{x}{|x|}$ is perpendicular to $S^2$. In general, the gradient is
perpendicular to the level sets.
\end{remark}

This intuition is made rigorous in the following.

\begin{proposition}[constrained optimization problems]
\label{prop:constrained_optimization}
Suppose $f : U \to \mathbb{R}$, $g : U \to \mathbb{R}^k$ are $C^1$ and $U \subseteq \mathbb{R}^n$
is open, and write $g = (g_1,\dots,g_k)$. If $f|_{g^{-1}\{0\}}$ has a local extremum at
$x_0 \in g^{-1}\{0\}$, then
$$\{\nabla f(x_0), \nabla g_1(x_0), \dots, \nabla g_k(x_0)\} \text{ are linearly dependent.}$$
\end{proposition}

For $k = 1$, this corresponds to our intuition above of $\nabla f(x_0)$ \textbf{parallel to}
$\nabla g(x_0)$.

\textit{In practice}, we assume $(\nabla g_1, \dots, \nabla g_k)$ nowhere linearly dependent,
since otherwise $\lambda_0 = 0$ is possible, which does not give us anything. We then eliminate
the perpendicular part of $\nabla f(x_0)$ by forming the \newterm{Lagrangian}
(\germanterm{Lagrange-Funktion}):
$$
\begin{aligned}
L : U \times \mathbb{R}^k &\to \mathbb{R} \\
(x,\lambda) &\mapsto f(x) - \lambda\cdot g(x) = f(x) - \lambda_1g_1(x) - \dots - \lambda_kg_k(x)
\end{aligned}
$$
We then find local extrema by forcing $\nabla L(x,\lambda) = 0$, i.e.
$$\underbrace{\frac{\partial L}{\partial x_i} = 0 \ \ \forall i = 1,\dots,n}_{\nabla f(x) - \lambda\cdot\nabla g(x) = 0} \qquad \text{and} \qquad \underbrace{\frac{\partial L}{\partial \lambda_j} = 0 \ \ \forall j = 1,\dots,k}_{g_j(x) = 0}$$

\begin{exercise}
\label{ex:xyz_on_ball}
Find the local extrema of
$$f : K \to \mathbb{R}, \quad (x,y,z) \mapsto xyz, \qquad K = \{(x,y,z) \in \mathbb{R}^3 : x^2+y^2+z^2 \leq 1\}.$$
\end{exercise}

\begin{exercisesolution}[\cref{ex:xyz_on_ball}]
We first find the \textbf{interior critical points}:
$$\nabla f(x,y,z) = (yz,\ xz,\ xy) \overset{!}{=} 0 \implies x = y = 0 \ \text{ or } \ x = z = 0 \ \text{ or } \ y = z = 0.$$
Then we look for local extrema \textbf{on the boundary}. We need to find the local extrema of
$f|_{g^{-1}\{0\}}$, where $g : \mathbb{R}^3\to\mathbb{R}$ is given by $g(x) = |x|^2 - 1$.

The Lagrangian is
$$L(x,y,z,\lambda) = f(x,y,z) - \lambda g(x,y,z) = xyz - \lambda(x^2+y^2+z^2-1),$$
giving the system of equations
$$
\begin{aligned}
\frac{\partial L}{\partial x} &= yz - 2\lambda x = 0, & \frac{\partial L}{\partial \lambda} &= 1 - (x^2+y^2+z^2) = 0. \\
\frac{\partial L}{\partial y} &= xz - 2\lambda y = 0, \\
\frac{\partial L}{\partial z} &= xy - 2\lambda z = 0,
\end{aligned}
$$
We then have
\begin{equation*}
2\lambda = \frac{yz}{x} = \frac{xz}{y} = \frac{xy}{z} \tag{$*$}
\end{equation*}
$$\implies \underbrace{y^2z^2 = x^2z^2}_{y^2 = x^2} = \underbrace{x^2y^2}_{y^2 = z^2} \qquad \text{(if any coordinate is 0, so are the others)}$$
$$\implies x^2 = y^2 = z^2.$$
Then $x^2+y^2+z^2 = 1 \implies 3x^2 = 1 \implies x = \pm\frac{1}{\sqrt{3}}$, giving
\textbf{8 solutions}:
$$(x,y,z) \in \left\{\tfrac{1}{\sqrt{3}}(\alpha,\beta,\gamma) : \alpha,\beta,\gamma \in \{\pm 1\}\right\}$$
plus the solutions $\{(x_1,x_2,x_3) \in K : x_i = x_j = 0,\ i \neq j\}$.

\textbf{To justify division by $z$, $x$ or $y$ in $(*)$:} if $z = 0$, then
$-\lambda x = -\lambda y = xy = 0$, giving $x = 0$ or $y = 0$, and $\lambda = 0$, since
$x = y = z = 0$ does not satisfy $x^2+y^2+z^2 = 1$. Then, since $\lambda = 0$,
$\nabla f(x,y,z) = 0$ and we have already found that point in the first step. Similarly for
$y = 0$ or $x = 0$.
\end{exercisesolution}

\begin{remark}
Corsin writes $g(x) = |x| - 1$ for the constraint above, but the Lagrangian uses
$x^2+y^2+z^2-1$, i.e.\ $|x|^2-1$. Both cut out the same sphere; the squared version is the one
actually differentiated, so it is used throughout. See \texttt{OQ-12}.
\end{remark}

\begin{figurestub}{FIG-W05-03}
A blue sphere with the three green coordinate axes through the centre; the eight points
$\tfrac{1}{\sqrt3}(\pm1,\pm1,\pm1)$ marked in red, connected by red dotted lines forming an
inscribed cube.
\end{figurestub}

\section{The Hessian test}
\label{sec:hessian_test}

The \newterm{Hessian} of $f \in C^2(U,\mathbb{R})$, $U \subseteq \mathbb{R}^n$, is given by
$$\Hess f(x) = (\partial_i\partial_j f)_{i,j=1,\dots,n}$$
with respect to the standard basis of $\mathbb{R}^n$. By \newterm{Schwarz's lemma}
(\germanterm{Satz von Schwarz}), $\Hess f(x)$ is symmetric.

\begin{definition}[sign of a symmetric matrix]
A symmetric matrix $A \in \mathbb{R}^{n\times n}$ is:
\begin{itemize}
  \item \newterm{positive definite} (\germanterm{positiv definit}) if $v\transp Av > 0$ for all
    $v \in \mathbb{R}^n\setminus\{0\}$;
  \item \newterm{negative definite} (\germanterm{negativ definit}) if $v\transp Av < 0$ for all
    $v \in \mathbb{R}^n\setminus\{0\}$;
  \item \textbf{indefinite} if there exist $v, u \in \mathbb{R}^n\setminus\{0\}$ such that
    $v\transp Av > 0$ and $u\transp Au < 0$;
  \item \textbf{degenerate} if there exists $v \in \mathbb{R}^n\setminus\{0\}$ such that
    $v\transp Av = 0$.
\end{itemize}
\end{definition}

\begin{theorem}[the Hessian test]
\label{thm:hessian_test}
If $f \in C^2(U,\mathbb{R})$ has a critical point $x_0 \in \operatorname{int}(U)$, then:
\begin{itemize}
  \item $\Hess f(x_0)$ \textbf{positive definite} $\Rightarrow$ local \textbf{min} at $x_0$;
  \item $\Hess f(x_0)$ \textbf{negative definite} $\Rightarrow$ local \textbf{max} at $x_0$;
  \item $\Hess f(x_0)$ \textbf{indefinite and not degenerate} $\Rightarrow$ \newterm{saddle point}
    (\germanterm{Sattelpunkt}) at $x_0$.
\end{itemize}
\end{theorem}

\begin{remark}[mnemonic]
positive $= \smile \dots$ min; negative $= \frown \dots$ max; indefinite $=$ saddle.
\end{remark}

\begin{remark}
The regularity in \cref{thm:hessian_test} is stated by Corsin as $C^3$; the test needs only
$C^2$ (which is what the Hessian's definition assumes). Relaxed to $C^2$ above. See
\texttt{OQ-13}.
\end{remark}

\subsection{How do I determine positive/negative definiteness?}

Let $A \in \mathbb{R}^{n\times n}$ be symmetric with eigenvalues $\lambda_1,\dots,\lambda_n$.
Note: $\det(A) = \lambda_1\cdots\lambda_n$ and $\Tr(A) = \lambda_1+\dots+\lambda_n$.

\textbf{For $n = 2$:}
\begin{itemize}
  \item $\det(A) = \lambda_1\lambda_2 < 0 \implies$ \textbf{indefinite}
  \item $\det(A) > 0$, $\Tr(A) > 0 \implies$ \textbf{positive definite}
  \item $\det(A) > 0$, $\Tr(A) < 0 \implies$ \textbf{negative definite}
\end{itemize}

\textbf{For $n \geq 3$:}
\begin{itemize}
  \item Solve the characteristic equation $\det(A - \lambda\id) = 0$ for $\lambda$ and check all
    eigenvalues.
  \item \newterm{Hurwitz criterion} (\germanterm{Hurwitz-Kriterium}): for
    $$A = \begin{pmatrix} a_{11} & a_{12} & \cdots & a_{1n} \\ a_{21} & \ddots & & \vdots \\ \vdots & & \ddots & \\ a_{n1} & \cdots & & a_{nn}\end{pmatrix}$$
    define the leading principal minors
    $$A_1 = \det(a_{11}), \quad A_2 = \det\begin{pmatrix} a_{11} & a_{12} \\ a_{21} & a_{22}\end{pmatrix}, \quad \dots, \quad A_n = \det(A).$$
    Then:
    \begin{itemize}
      \item $A_1 > 0$, $A_2 > 0$, $A_3 > 0$, \ldots\ $\iff$ \textbf{positive definite}
      \item $A_1 < 0$, $A_2 > 0$, $A_3 < 0$, \ldots\ $\iff$ \textbf{negative definite}
      \item \textbf{indefinite} if neither $A_1 \geq 0, A_2 \geq 0, \dots$ nor
        $A_1 \leq 0, A_2 \geq 0, A_3 \leq 0, \dots$
    \end{itemize}
\end{itemize}

\begin{remark}
Corsin writes the characteristic equation as ``$A - \lambda\id = 0$'', omitting the determinant.
Corrected above. See \texttt{OQ-14}.
\end{remark}

\begin{example}
\begin{center}
\begin{tabular}{lll}
\textbf{Matrix} & \textbf{Computation} & \textbf{Conclusion} \\
\hline
$\begin{pmatrix}2&2\\2&5\end{pmatrix}$ & $\det = 6 > 0$, $\Tr = 7 > 0$ & positive definite \\
$\begin{pmatrix}-4&2\\2&-4\end{pmatrix}$ & $\det = 12 > 0$, $\Tr = -8 < 0$ & negative definite \\
$\begin{pmatrix}4&2\\2&-4\end{pmatrix}$ & $\det = -20 < 0$ & indefinite \\
$\begin{pmatrix}-5&0&0\\0&-4&-2\\0&-2&-4\end{pmatrix}$ & $A_1 = -5 < 0$, $A_2 = 20 > 0$, $A_3 = -60 < 0$ & negative definite \\
$\begin{pmatrix}-1&0&-4\\0&2&-2\\-4&-2&22\end{pmatrix}$ & $A_1 = -1 < 0$, $A_2 = -2 < 0$, $A_3 = -72 < 0$ & indefinite \\
\end{tabular}
\end{center}
\end{example}

\begin{exercise}
\label{ex:cubic_saddle_family}
Find the critical values of
$$f_\alpha : \mathbb{R}^2 \to \mathbb{R}, \qquad (x,y) \mapsto x^3 - y^3 + 3\alpha xy$$
for $\alpha \in \mathbb{R}\setminus\{0\}$ and determine min/max/saddle points.
\end{exercise}

\begin{exercisesolution}[\cref{ex:cubic_saddle_family}]
$$
\begin{aligned}
\frac{\partial f_\alpha}{\partial x} &= 3x^2 + 3\alpha y \overset{!}{=} 0 &&\implies y = -\frac{x^2}{\alpha} \\
\frac{\partial f_\alpha}{\partial y} &= -3y^2 + 3\alpha x \overset{!}{=} 0 &&\implies -\frac{x^4}{\alpha^2} + \alpha x = 0
\end{aligned}
$$
\begin{enumerate}[label=\textbf{\arabic*.}]
  \item $x = y = 0$
  \item $x \neq 0 \implies x^3 = \alpha^3 \implies x = \alpha$, $y = -\alpha$
\end{enumerate}
So we have critical points $(0,0)$ and $(\alpha, -\alpha)$.

Compute the Hessian:
$$\begin{pmatrix} \partial_x^2 f & \partial_x\partial_y f \\ \partial_x\partial_y f & \partial_y^2 f\end{pmatrix} = \begin{pmatrix} 6x & 3\alpha \\ 3\alpha & -6y \end{pmatrix} = \Hess f(x,y)$$
$$\Hess f(0,0) = \begin{pmatrix} 0 & 3\alpha \\ 3\alpha & 0\end{pmatrix}, \qquad \det = -9\alpha^2 < 0 \implies \text{indefinite}$$
$$\Hess f(\alpha,-\alpha) = \begin{pmatrix} 6\alpha & 3\alpha \\ 3\alpha & 6\alpha\end{pmatrix} = 3\alpha\begin{pmatrix} 2 & 1 \\ 1 & 2\end{pmatrix}, \qquad \det = 27\alpha^2 > 0, \quad \Tr = 12\alpha$$
For $\alpha > 0$ this is positive definite; for $\alpha < 0$ it is negative definite. So:
\begin{itemize}
  \item $(0,0) \longrightarrow$ \textbf{saddle point}
  \item $(\alpha,-\alpha) \longrightarrow$ \textbf{local minimum} for $\alpha > 0$; \textbf{local
    maximum} for $\alpha < 0$
\end{itemize}
\end{exercisesolution}
